\documentclass[12pt]{article}

\usepackage{listings}
\usepackage{enumitem}
\usepackage{fancyhdr}
\usepackage{lastpage}
\usepackage{ragged2e}
\usepackage{lipsum}
\usepackage{newtxtext,newtxmath}
\usepackage{graphicx}
\usepackage{titlesec}
\usepackage{titling}
\usepackage[top=2.5cm, bottom=2.5cm, left=3cm, right=3cm]{geometry}

\setlength{\emergencystretch}{3em}
\AtBeginDocument{\RaggedRight}

\setlength{\headheight}{14.5pt}
\pagestyle{fancy}
\fancyhf{}
\fancyhead[L]{\small\textit{Group 15 -- Lab Report}}
\fancyhead[R]{\small\textit{\thepage / \pageref{LastPage}}}
\fancyfoot[C]{\small\thepage}

\renewcommand{\headrulewidth}{1.0pt}
\setlength{\droptitle}{-3em}

\title{\fontsize{16pt}{18pt}\selectfont\textbf{Lab 2 - H-Bridge with DC Motor}}
\author{\fontsize{11pt}{13pt}\selectfont
  Yuwei Zhao\textsuperscript{1},
  Haohao Qi\textsuperscript{2},
  Yilin Zhang\textsuperscript{3} \\
  \small\textsuperscript{1,2,3}School of Computer Science and Technology, University of Science \\
  \small\textsuperscript{1}Yuwei Zhao (Experimental design \& initial draft) \\
  \small\textsuperscript{2}Haohao Qi (Data acquisition \& visualization) \\
  \small\textsuperscript{3}Yilin Zhang (Result discussion \& paper revision and polishment)
}
\date{\fontsize{11pt}{13pt}\selectfont\today}

\titleformat{\section}[hang]
{\normalfont\fontsize{14}{16}\bfseries}
{\thesection}{0.5em}{}
\titlespacing{\section}{0pt}{2.5ex plus 1ex minus .2ex}{1.5ex plus .2ex}

\titlespacing*{\section}
  {0pt}
  {2ex}
  {1.5ex}

\begin{document}

\maketitle
\thispagestyle{fancy}
\vspace*{-0.5em}

\begin{abstract} 
\noindent Motors are widely used in robotics, industrial machinery and automation, and many other industries. 
By understanding the schematic diagram of a specific motor givin in the lab, 
we build a circuit based on $Arduino Kona328$ and $chip-L293D$ to run the motor.
We demonstrate adjusting speed, and then reverse its rotation direction. 
Additionally, since light-dependent registers (LDRs) can be used to change speed and direction, 
a light focused car motor design is presented in this lab report.
\end{abstract}

\bigskip

\section*{Introduction}
This experiment aims on discovering the basic working principle of a simple motor. 
Then a theory of experiments would be 
An Arduino programmed circuit is built to run the motor. 
\begin{itemize}
  \item Set up motor to be running;
  \item Change motor to a certain speed;
  \item Run motor forward and reverse.
\end{itemize}
The motor is placed in a breadboard, and powered by a voltage up to 5V.
Motor's speed is adjusted by changing voltage, with a 

Give a brief description for the aims of the laboratory and a summary of the results collected.
Please write this as constant prose in a way that a fellow student,
who does not know anything about the experiment,
can understand the approach.

\section*{Theory and Technical Background}
Give a brief description of the background theory for the experiment being tackled.
Please cite appropriate work.
Feel free to add figures into this section, where
each figure and each equation may reduce the number of words you use.

\section*{Method}
Describe the technical procedures to collect the results.
Add figures into this section, where
each figure and each equation may reduce the number of words you use.

\section*{Results}
Add all appropriate figures and analysis.
No formal word count for this section, however
please keep text as concise as possible.

\section*{Conclusions}
Please give the conclusions from the characterisations including a summary of the results collected.

\end{document}